% Source: Wald GR, physical PDF pp. 54--60
%         (printed pp. 41--47).
% Coverage: Section 3.3, Geodesics; equations (3.3.1)--(3.3.18).
% Figures/tables/footnotes: Figures 3.4--3.5; no tables or footnotes.
% Status: source-reviewed and compile-clean.
% Uncertainties: none.

\subsection{Geodesics}\label{sec:3.3}

Intuitively, geodesics are lines that \lq\lq curve as little as possible\rq\rq; they are the
\lq\lq straightest possible lines\rq\rq\ one can draw in a curved geometry. Given a derivative
operator, $\nabla_a$, we define a \emph{geodesic} to be a curve whose tangent vector is parallel
propagated along itself, i.e. a curve whose tangent, $T^a$, satisfies the equation
\begin{equation}
  T^a\nabla_aT^b=0.
  \tag{3.3.1}\label{eq:3.3.1}
\end{equation}
Actually, in order to satisfy the intuitive requirement that the curve be \lq\lq as straight
as possible,\rq\rq\ one might require only that the tangent vector to the curve point in the
same direction as itself when parallel propagated, and not demand that it maintain
the same length. This would yield the weaker condition,
\begin{equation}
  T^a\nabla_aT^b=\alpha T^b,
  \tag{3.3.2}\label{eq:3.3.2}
\end{equation}
where $\alpha$ is an arbitrary function on the curve. However, it is easy to show that given
a curve which satisfies Eq.~\eqref{eq:3.3.2} we can always reparameterize it so that it
satisfies Eq.~\eqref{eq:3.3.1} (see problem 5). Thus, there is no true loss of generality
in considering only curves which satisfy Eq.~\eqref{eq:3.3.1}. A parameterization which
yields Eq.~\eqref{eq:3.3.1} is called an \emph{affine parameterization}. Thus, our definition of
a geodesic requires it to be affinely parameterized. (Some other references apply the
term geodesic to any curve satisfying Eq.~(3.3.2).)

We can get some insight into the nature of the geodesic equation by writing out
the components of this equation in a coordinate basis. In a coordinate system $\psi$, the
geodesic is mapped into a curve $x^\mu(t)$ in $\mathbb R^n$. By Eq.~\eqref{eq:3.1.19}, the components,
$T^\mu$, of $T^a$ in this coordinate basis satisfy
\begin{equation}
  \frac{dT^\mu}{dt}+\sum_{\sigma,\nu}\Gamma^\mu{}_{\sigma\nu}T^\sigma T^\nu=0.
  \tag{3.3.3}\label{eq:3.3.3}
\end{equation}
However, by Eq.~\eqref{eq:2.2.12}, the components $T^\mu$ are simply
\begin{equation}
  T^\mu=\frac{dx^\mu}{dt}.
  \tag{3.3.4}\label{eq:3.3.4}
\end{equation}
Thus, in a coordinate basis, the geodesic equation becomes
\begin{equation}
  \frac{d^2x^\mu}{dt^2}
  +\sum_{\sigma,\nu}\Gamma^\mu{}_{\sigma\nu}
  \frac{dx^\sigma}{dt}\frac{dx^\nu}{dt}=0.
  \tag{3.3.5}\label{eq:3.3.5}
\end{equation}

Equation~\eqref{eq:3.3.5} is a coupled system of $n$ second order ordinary differential
equations for the $n$ functions $x^\mu(t)$. From the theory of ordinary differential equations, it is known that a unique solution of Eq.~\eqref{eq:3.3.5} always exists for any
given initial value of $x^\mu$ and $dx^\mu/dt$. This means that given $p\in M$ and any tangent
vector, $T^a\in T_pM$, there always exists a unique geodesic through $p$ with tangent $T^a$.
Notice that the solutions of the equations of motion in ordinary mechanics share this
property: Given an initial position and velocity, a unique solution exists.

Indeed, the existence and uniqueness of geodesics allow us to use them to construct coordinate systems that are very convenient for some computational purposes.
Let $p\in M$. We define a map, called the \emph{exponential map}, from the tangent space
$T_pM$ to $M$ by mapping each $T^a\in T_pM$ into the point in $M$ lying at unit affine parameter
from $p$ along the geodesic through $p$ with tangent $T^a$. For large $T^a$ one might
encounter a singularity before the affine parameter $t=1$ is reached. Also geodesics
may cross, thereby making the exponential map fail to be one-to-one. However, one
can show that there always exists a (sufficiently small) neighborhood of the origin
of $T_pM$ on which the exponential map is defined and is one-to-one (see, e.g., Bishop
and Crittenden 1964). Since $T_pM$ is an $n$-dimensional vector space, we may identify
it with $\mathbb R^n$, and hence use the exponential map to give us a coordinate system, called
\emph{Riemannian normal coordinates} at $p$. These coordinates have the property that all
geodesics through $p$ get mapped into straight lines through the origin of $\mathbb R^n$. From
Eq.~\eqref{eq:3.3.5} it follows that in this coordinate system the Christoffel symbol
components $\Gamma^\mu{}_{\sigma\nu}$ vanish at $p$. This fact makes Riemannian normal coordinates
particularly useful if one is performing calculations at a given point.

In the case where the derivative operator $\nabla_a$ arises from a metric $g_{ab}$ a second type
of coordinate system, called \emph{Gaussian normal coordinates}, or \emph{synchronous coordinates}, often is useful for calculations in situations where one is given a \emph{hypersurface}
$S$, i.e., an $(n-1)$-dimensional embedded submanifold of the $n$-dimensional manifold $M$ (see appendix B). At each point $p\in S$, the tangent space $T_pS$ of the manifold
$S$ can be naturally viewed as an $(n-1)$-dimensional subspace of the tangent space
$T_pM$ of $M$. Thus, there will be a vector $n^a\in T_pM$, unique up to scaling, which is
orthogonal (with respect to the metric $g_{ab}$) to all vectors in $T_pS$. This vector, $n^a$, is said
to be \emph{normal} to $S$. In the case of a Riemannian metric, $n^a$ cannot lie in $T_pS$; in the case
of a metric of indefinite signature, $n^a$ could be a null vector, $g_{ab}n^an^b=0$, in which
case it does lie in $T_pS$ and $S$ is said to be a \emph{null hypersurface} at point $p$. If $S$ is nowhere
null, we may normalize $n^a$ by the condition $g_{ab}n^an^b=\pm1$. Gaussian normal coordinates may be defined for any non-null hypersurface as follows (see Fig.~\ref{fig:3.4}). For
each $p\in S$ we construct the unique geodesic through $p$ with tangent $n^a$. We choose
arbitrary coordinates $(x^1,\ldots,x^{n-1})$ on a (portion of) $S$ and label each point in a
neighborhood of (that portion of) $S$ by the parameter $t$ along the geodesic on which
it lies and the coordinates $x^1,\ldots,x^{n-1}$ of the point $p\in S$ from which the geodesic
emanated. The geodesics emanating from $S$ may eventually cross or run into singularities, but in a (sufficiently small) neighborhood of each $p\in S$, the map
$q\mapsto(x^1,\ldots,x^{n-1},t)$ defines the chart we wished to construct.

Gaussian normal coordinates satisfy the important property that the geodesics
remain orthogonal to all the hypersurfaces $S_t$ defined by the equation $t=\text{constant}$.
This is true by construction for the hypersurface $S_0=S$. To show that this remains
true for all $S_t$ in the domain of validity of the construction, it suffices to show that
the geodesic tangent field $n^a$ remains orthogonal to all of the coordinate basis fields

% Source: Wald GR, physical PDF p. 56 (printed p. 43).
% Coverage: Figure 3.4.
% Figures/tables/footnotes: Figure only; no tables or footnotes.
% Status: source-reviewed and compile-clean.
% Uncertainties: none.
\begingroup
\renewcommand{\thefigure}{3.4}
\begin{figure}[H]
  \centering
  \begin{tikzpicture}[scale=.9,>=Stealth]
    \foreach \y in {0,1.65} {
      \draw[thick] (-2.15,\y) .. controls (-.80,\y+.32) and (.75,\y+.22) .. (2.10,\y-.08)
        -- (1.55,\y-.78) .. controls (.15,\y-.48) and (-.90,\y-.63) .. (-2.15,\y-.34) -- cycle;
    }
    % Geodesics orthogonal to S; the leftmost pair crosses above S_t, which is
    % where Gaussian normal coordinates break down.
    \def\geodesics{%
      \draw[thick] (-1.15,.05) .. controls (-1.20,1.20) and (-1.10,2.10) .. (-.20,3.00);
      \draw[thick] (-.35,.05) .. controls (-.42,1.20) and (-.48,2.10) .. (-.75,3.00);
      \draw[thick] (.95,.05) .. controls (.98,1.20) and (1.10,2.10) .. (1.35,3.00);%
    }
    \begin{scope}
      \clip (-2.6,-1.1) rectangle (2.6,1.12) (-2.6,1.82) rectangle (2.6,3.2);
      \geodesics
    \end{scope}
    \begin{scope}[dashed]
      \clip (-2.6,1.12) rectangle (2.6,1.82);
      \geodesics
    \end{scope}
    \node at (-1.62,1.42) {$S_t$};
    \node at (-1.62,-.20) {$S$};
    \fill (-.35,.05) circle (1.8pt) node[below=1pt] {$p$};
  \end{tikzpicture}
  \caption{The construction of Gaussian normal coordinates starting from the
  hypersurface $S$. The geodesics orthogonal to $S$ eventually may cross, but until they
  do, Gaussian normal coordinates are well defined and the surfaces, $S_t$, of constant $t$
  remain orthogonal to the geodesics.}
  \label{fig:3.4}
\end{figure}
\endgroup


$X^a_1,\ldots,X^a_{n-1}$ which generate the tangent space to $S_t$. Denoting by $X^a$ any one of
these fields, we have
\begin{equation}
  \begin{aligned}
    n^b\nabla_b(n_aX^a)
    &=n_an^b\nabla_bX^a\\
    &=n_aX^b\nabla_bn^a\\
    &=\frac{1}{2}X^b\nabla_b(n^an_a)\\
    &=0,
  \end{aligned}
  \tag{3.3.6}\label{eq:3.3.6}
\end{equation}
where the first equality comes from the geodesic equation, the second from the fact
that $n^a$ and $X^b$---being elements of a coordinate basis on $M$---commute, the third
follows directly from the Leibniz rule, and the last follows from the fact that the
normalization condition $n^an_a=\pm1$ on $S$ is preserved by parallel transport so that
$n^an_a$ is constant on $M$. Since $n_aX^a=0$ on $S$, Eq.~\eqref{eq:3.3.6} shows that this
condition is preserved off of $S$, which yields the desired result.

A further property of geodesics of a derivative operator arising from a metric is
that they extremize the length of curves connecting given points as measured by the
metric. For a smooth (or merely differentiable) curve $C$ on a manifold $M$ with
Riemannian metric $g_{ab}$, the \emph{length}, $l$, of $C$ is defined by
\begin{equation}
  l=\int(g_{ab}T^aT^b)^{1/2}\,dt,
  \tag{3.3.7}\label{eq:3.3.7}
\end{equation}
where $T^a$ is the tangent to $C$ and $t$ is the curve parameter. For a metric of Lorentz
signature $-++\cdots+$, a curve is said to be \emph{timelike} if the norm of its tangent is
everywhere negative, $g_{ab}T^aT^b<0$; it is said to be \emph{null} if $g_{ab}T^aT^b=0$; it is said to
be \emph{spacelike} if $g_{ab}T^aT^b>0$. For spacelike curves, the length may again be defined
by Eq.~\eqref{eq:3.3.7}; for null curves the length is zero; for timelike curves, we change
the sign in the square root and use the term \emph{proper time}, $\tau$, rather than length,
\begin{equation}
  \tau=\int(-g_{ab}T^aT^b)^{1/2}\,dt.
  \tag{3.3.8}\label{eq:3.3.8}
\end{equation}
The length of curves which change from timelike to spacelike is not defined. Note
that since a geodesic tangent is parallel transported and thus its norm is constant,
geodesics in a Lorentz manifold cannot change from timelike to spacelike or null.
Note also that the length (or proper time) of a curve does not depend on the way in
which the curve is parameterized. If we define a new reparameterization $s=s(t)$, the
new tangent will be $S^a=(dt/ds)T^a$ and the new length will be
\begin{equation}
  l'=\int[g_{ab}S^aS^b]^{1/2}\,ds
  =\int[g_{ab}T^aT^b]^{1/2}\frac{dt}{ds}\,ds=l.
  \tag{3.3.9}\label{eq:3.3.9}
\end{equation}

We wish now to derive the condition on a curve which makes it extremize the
length between its endpoints, i.e., wish to find those curves whose length does not
change to first order under an arbitrary smooth deformation which keeps the endpoints fixed. We will perform the calculation of the change in length of a curve under
a small deformation by choosing a chart and working in $\mathbb R^n$. (In chapter 9 we shall
repeat this calculation without introducing coordinates and will calculate also the
second variation of arc length.) For definiteness, we consider a spacelike curve.
Writing Eq.~\eqref{eq:3.3.7} in the coordinate basis, we have
\begin{equation}
  l=\int_a^b\left[\sum_{\mu,\nu}g_{\mu\nu}
  \frac{dx^\mu}{dt}\frac{dx^\nu}{dt}\right]^{1/2}\,dt,
  \tag{3.3.10}\label{eq:3.3.10}
\end{equation}
where $C(a)=p$ and $C(b)=q$ are the endpoints of the curve. The extremization
problem for $l$ is mathematically identical to the extremization problem for the action
in Lagrangian mechanics. The variation in $l$ is
\begin{equation}
  \begin{aligned}
    \delta l={}&\int_a^b\left[\sum_{\mu,\nu}g_{\mu\nu}
    \frac{dx^\mu}{dt}\frac{dx^\nu}{dt}\right]^{-1/2}
    \sum_{\alpha,\beta}\biggl\{
    g_{\alpha\beta}\frac{dx^\alpha}{dt}\frac{d(\delta x^\beta)}{dt}\\
    &\qquad\qquad+\frac{1}{2}\sum_\sigma
    \frac{\partial g_{\alpha\beta}}{\partial x^\sigma}
    \delta x^\sigma\frac{dx^\alpha}{dt}\frac{dx^\beta}{dt}
    \biggr\}\,dt.
  \end{aligned}
  \tag{3.3.11}\label{eq:3.3.11}
\end{equation}
Without loss of generality (since length is parameterization independent), we may
assume that the original curve was parameterized so that
\begin{equation*}
  g_{ab}T^aT^b=1
  =\sum_{\mu,\nu}g_{\mu\nu}\frac{dx^\mu}{dt}\frac{dx^\nu}{dt}.
\end{equation*}
With this choice of parameterization, the extremization condition is
\begin{equation}
  \begin{aligned}
    0={}&\int_a^b\sum_{\alpha,\beta}\biggl\{
    g_{\alpha\beta}\frac{dx^\alpha}{dt}\frac{d(\delta x^\beta)}{dt}
    +\frac{1}{2}\sum_\lambda
    \frac{\partial g_{\alpha\lambda}}{\partial x^\beta}
    \frac{dx^\alpha}{dt}\frac{dx^\lambda}{dt}\delta x^\beta
    \biggr\}\,dt\\
    ={}&\int_a^b\sum_{\alpha,\beta}\biggl\{
    -\frac{d}{dt}\left(g_{\alpha\beta}\frac{dx^\alpha}{dt}\right)
    +\frac{1}{2}\sum_\lambda
    \frac{\partial g_{\alpha\lambda}}{\partial x^\beta}
    \frac{dx^\alpha}{dt}\frac{dx^\lambda}{dt}
    \biggr\}\delta x^\beta\,dt.
  \end{aligned}
  \tag{3.3.12}\label{eq:3.3.12}
\end{equation}
(No boundary terms occur in the integration by parts since $\delta x^\beta$ vanishes at the
endpoints.) Equation~\eqref{eq:3.3.12} will hold for arbitrary $\delta x^\beta$ if and only if
\begin{equation}
  -\sum_\alpha g_{\alpha\beta}\frac{d^2x^\alpha}{dt^2}
  -\sum_{\alpha,\lambda}\frac{\partial g_{\alpha\beta}}{\partial x^\lambda}
  \frac{dx^\lambda}{dt}\frac{dx^\alpha}{dt}
  +\frac{1}{2}\sum_{\alpha,\lambda}
  \frac{\partial g_{\alpha\lambda}}{\partial x^\beta}
  \frac{dx^\alpha}{dt}\frac{dx^\lambda}{dt}=0.
  \tag{3.3.13}\label{eq:3.3.13}
\end{equation}
Using our formula for $\Gamma^\sigma{}_{\alpha\lambda}$, Eq.~\eqref{eq:3.1.30}, we see that Eq.~\eqref{eq:3.3.13} is just
the geodesic equation~\eqref{eq:3.3.5}. (Had we not chosen the parameterization
$g_{ab}T^aT^b=1$ for $C$, we would have obtained the geodesic equation in the form
(3.3.2).) Thus, a curve extremizes the length between its endpoints if and only if it
is a geodesic.

An identical derivation shows that the curves which extremize proper time between two points are precisely the timelike geodesics. These derivations also show
that the geodesic equation (with affine parameterization) can be obtained from
variation of the Lagrangian,
\begin{equation}
  L=\sum_{\mu,\nu}g_{\mu\nu}\frac{dx^\mu}{dt}\frac{dx^\nu}{dt}.
  \tag{3.3.14}\label{eq:3.3.14}
\end{equation}
In many cases, the most efficient way of computing the Christoffel symbol
$\Gamma^\mu{}_{\sigma\nu}$ in a given coordinate basis is to start with the Lagrangian, Eq.~\eqref{eq:3.3.14}, write
down the corresponding Euler-Lagrange equations, and read off $\Gamma^\mu{}_{\sigma\nu}$ by comparison
with Eq.~\eqref{eq:3.3.5}.

On a manifold with a Riemannian metric, one can always find curves of arbitrarily
long length connecting any two points. However, the length will be bounded from
below, and the curve of shortest length connecting two points (assuming the lower
bound in length is attained) is necessarily an extremum of length and thus a geodesic.
Thus, the shortest path between two points is always a \lq\lq straightest possible path.\rq\rq\
However, a given geodesic connecting two points need not be the shortest path
between them. For a manifold with a Lorentz metric, given two points that can be
connected by a timelike curve, one can always find timelike curves of arbitrarily
small proper time connecting the points (see Fig.~9.5 of chapter 9). In some spacetimes the proper time of timelike curves connecting the two given points need not
be bounded from above; but if a curve of greatest proper time exists, it must be a
timelike geodesic. Again, however, a given geodesic need not maximize the proper
time between two points. In chapter 9 we shall introduce the notion of conjugate
points along a geodesic and will show that their absence is the necessary and
sufficient condition for a geodesic to be a local maximum of proper time (or, in the
Riemannian case, a local minimum of length) between two points.

Our final task is to derive the geodesic deviation equation, the equation which
relates the tendency of geodesics to accelerate toward or away from each other to the
curvature of the manifold. This gives another characterization of curvature, and it
also plays an important role in motivating Einstein's equation (see section 4.3) and
in arguments relating to the singularity theorems (see chapter 9).

Let $\gamma_s(t)$ denote a smooth one-parameter family of geodesics (see Fig.~\ref{fig:3.5}), that
is for each $s\in\mathbb R$, the curve $\gamma_s$ is a geodesic (parameterized by affine parameter $t$);

% Source: Wald GR, physical PDF p. 59 (printed p. 46).
% Coverage: Figure 3.5.
% Figures/tables/footnotes: Figure only; no tables or footnotes.
% Status: source-reviewed and compile-clean.
% Uncertainties: none.
\begingroup
\renewcommand{\thefigure}{3.5}
\begin{figure}[H]
  \centering
  \begin{tikzpicture}[scale=.9,>=Stealth]
    \foreach \x/\bend in {-1.0/-.26,-.45/-.10,.15/.08,.74/.22} {
      \draw[thick] (\x,-1.35) .. controls (\x-.05,-.35) and (\x+.10,.72) .. (\x+.35+\bend,2.05);
      \draw[->,thick] (\x,-1.42) -- ++(.18,.24);
    }
    \draw[->] (-1.55,-1.60) to[bend left=15] (-1.12,-1.20);
    \node[left=1pt] at (-1.58,-1.62) {$\gamma_s$};
    \coordinate (O) at (-.37,.08);
    \draw[->,very thick] (O) -- ++(-.08,.68);
    \node[fill=white,inner sep=1pt] at (-.68,.80) {$T^a$};
    \draw[->,very thick] (O) -- ++(.62,.02);
    \node[fill=white,inner sep=1pt] at (.30,-.22) {$X^a$};
  \end{tikzpicture}
  \caption{A one-parameter family of geodesics $\gamma_s$, with tangent $T^a$ and deviation
  vector $X^a$.}
  \label{fig:3.5}
\end{figure}
\endgroup


and the map $(t,s)\mapsto\gamma_s(t)$ is smooth, is one-to-one, and has smooth inverse. Let $\Sigma$
denote the two-dimensional submanifold spanned by the curves $\gamma_s(t)$. We may
choose $s$ and $t$ as coordinates of $\Sigma$. The vector field $T^a=(\partial/\partial t)^a$ is tangent to the
family of geodesics and, thus, satisfies
\begin{equation}
  T^a\nabla_aT^b=0.
  \tag{3.3.15}\label{eq:3.3.15}
\end{equation}
The vector field $X^a=(\partial/\partial s)^a$ represents the displacement to an infinitesimally
nearby geodesic, and is called the \emph{deviation vector}. There is \lq\lq gauge freedom\rq\rq\ in $X^a$
in the sense that $X^a$ changes by addition of a multiple of $T^a$ under a change of the
affine parameterization of the geodesics $\gamma_s(t)$, $t\mapsto t'=b(s)t+c(s)$ (see problem
5). It is worth noting that in the case where the geodesics arise from the derivative
operator associated with a metric $g_{ab}$, $X^a$ always can be chosen orthogonal to $T^a$.
Namely, by re-scaling $t$ by an $s$-dependent factor, we may ensure that $g_{ab}T^aT^b$
(which is automatically constant along each geodesic) does not vary with $s$. Since $X^a$
and $T^a$ are coordinate vector fields, they commute:
\begin{equation}
  T^b\nabla_bX^a=X^b\nabla_bT^a;
  \tag{3.3.16}\label{eq:3.3.16}
\end{equation}
so by the same calculation as Eq.~\eqref{eq:3.3.6}, we see that $X^aT_a$ is constant along
each geodesic. By further reparameterizing each $\gamma_s(t)$ by adding a constant (depending on $s$) to $t$, we may ensure that the curve $C(s)$ comprising the points with
$t=0$ is orthogonal to all the geodesics. Thus, with this affine parameterization of
$\gamma_s(t)$ we have $X_aT^a=0$ at $t=0$ and hence $X^aT_a=0$ everywhere.

The quantity $v^a=T^b\nabla_bX^a$ gives the rate of change along a geodesic of the
displacement to an infinitesimally nearby geodesic. Thus, we may interpret $v^a$ as the
relative velocity of an infinitesimally nearby geodesic. Similarly, we may interpret
\begin{equation}
  a^a=T^c\nabla_cv^a=T^c\nabla_c(T^b\nabla_bX^a)
  \tag{3.3.17}\label{eq:3.3.17}
\end{equation}
as the relative acceleration of an infinitesimally nearby geodesic in the family. We
now shall derive an equation which relates $a^a$ to the Riemann tensor. We have
\begin{equation}
  \begin{aligned}
    a^a
    &=T^c\nabla_c(T^b\nabla_bX^a)\\
    &=T^c\nabla_c(X^b\nabla_bT^a)\\
    &=(T^c\nabla_cX^b)(\nabla_bT^a)+X^bT^c\nabla_c\nabla_bT^a\\
    &=(X^c\nabla_cT^b)(\nabla_bT^a)+X^bT^c\nabla_b\nabla_cT^a
      -R_{cbd}{}^aX^bT^cT^d\\
    &=X^c\nabla_c(T^b\nabla_bT^a)-R_{cbd}{}^aX^bT^cT^d\\
    &=-R_{cbd}{}^aX^bT^cT^d.
  \end{aligned}
  \tag{3.3.18}\label{eq:3.3.18}
\end{equation}
Equation~\eqref{eq:3.3.18} is known as the \emph{geodesic deviation equation}. It yields the final
characterization of curvature which we sought: We have $a^a=0$ for all families of
geodesics if and only if $R_{abc}{}^d=0$. Thus, some geodesics will accelerate toward or
away from each other (and, in particular, some initially parallel geodesics---i.e.,
geodesics with $v^a=T^b\nabla_bX^a=0$ initially---will fail to remain parallel) if and only
if $R_{abc}{}^d\ne0$.
